% ============================================================
\section{Change of Basis: Geometric Interpretation}
% ============================================================

This section provides a geometric interpretation of similarity
transformation $A_{\text{new}} = P^{-1}AP$, then contrasts it with the
circuit transformation $Z_{\text{new}} = T^v Z (T^i)^{-1}$.

% ---- B.1 ----
\subsection{The Same Vector, Different Coordinates}
\label{sec:same-vector}

A fundamental insight from linear algebra: a vector is a geometric
object — an arrow in space — that exists independently of any coordinate
system. When we change basis, we are not moving the vector; we are only
changing how we describe it numerically.

\begin{center}
%\vspace{2mm}
\textit{The arrow stays fixed. Only the numbers change.}
%\vspace{2mm}
\end{center}

Consider $\mathbb{R}^2$:
\begin{itemize}
    \item Standard basis: $e_1 = \begin{bmatrix} 1 \\ 0 \end{bmatrix}$,
                         $e_2 = \begin{bmatrix} 0 \\ 1 \end{bmatrix}$
    \item New basis:     $e'_1 = \begin{bmatrix} 1 \\ 1 \end{bmatrix}$,
                        $e'_2 = \begin{bmatrix} -1 \\ 1 \end{bmatrix}$
\end{itemize}

A vector $\mathbf{x}$ pointing in some direction can be written as:
\[
    \mathbf{x} = 3e_1 + 2e_2 \quad\text{or}\quad
    \mathbf{x} = 2.5e'_1 + 0.5e'_2
\]

The arrow is identical. The coordinates $(3,2)$ vs $(2.5, 0.5)$
describe the same geometric object in different languages.

% ---- B.2 ----
\subsection{Change of Basis Matrix}
\label{sec:change-of-basis}

Let $v = \{e_1, e_2, \ldots, e_n\}$ be the old basis and $v' = \{e'_1,
e'_2, \ldots, e'_n\}$ be the new basis. The change-of-basis matrix
$M_{v' \to v}$ converts coordinates in the new basis to coordinates in
the old basis:

\[
    \mathbf{x}_v = M_{v' \to v}\,\mathbf{x}_{v'}
\]

\textbf{How to build $M$:} Each column $j$ of $M$ contains the
coordinates of the new basis vector $e'_j$ written in the old basis $v$.

\textbf{Example:} In $\mathbb{R}^2$ with $e'_1 = \begin{bmatrix} 1 \\ 1 \end{bmatrix}$
and $e'_2 = \begin{bmatrix} -1 \\ 1 \end{bmatrix}$ (expressed in standard basis):
\[
    M_{v' \to v} = \begin{bmatrix} 1 & -1 \\ 1 & 1 \end{bmatrix}
\]

To go the other direction (old to new), use the inverse:
\[
    \mathbf{x}_{v'} = M_{v' \to v}^{-1}\,\mathbf{x}_v
\]

% ---- B.3 ----
\subsection{Linear Transformation as a Physical Action}
\label{sec:linear-transformation}

Now consider a linear transformation $T$ — a physical operation like
rotation, scaling, or projection. Crucially, $T$ itself does not depend
on any basis; it is the same action regardless of how we choose to
describe vectors.

When we represent $T$ as a matrix, we are simply writing down the same
physical action in a particular coordinate language.

\textbf{How to find the matrix $[T]_v$ of transformation $T$ in basis $v$:}
Apply $T$ to each basis vector, then write the result in terms of the
same basis. These results become the columns of the matrix.

Specifically:
\begin{itemize}
    \item Column 1 of $[T]_v$: coordinates of $T(e_1)$ in basis $v$
    \item Column 2 of $[T]_v$: coordinates of $T(e_2)$ in basis $v$
    \item And so on...
\end{itemize}

\textbf{Example:} In $\mathbb{R}^2$, let $T$ be a $90^\circ$ rotation
counterclockwise. In the standard basis:
\[
    T(e_1) = T\!\begin{bmatrix} 1 \\ 0 \end{bmatrix}
           = \begin{bmatrix} 0 \\ 1 \end{bmatrix}
    \quad\Rightarrow\quad
    \text{column 1 } = \begin{bmatrix} 0 \\ 1 \end{bmatrix}
\]
\[
    T(e_2) = T\!\begin{bmatrix} 0 \\ 1 \end{bmatrix}
           = \begin{bmatrix} -1 \\ 0 \end{bmatrix}
    \quad\Rightarrow\quad
    \text{column 2 } = \begin{bmatrix} -1 \\ 0 \end{bmatrix}
\]
\[
    [T]_v = \begin{bmatrix} 0 & -1 \\ 1 & 0 \end{bmatrix}
\]

The matrix is different in a different basis, but the physical rotation
is the same.

% ---- B.4 ----
\subsection{Deriving $B = M^{-1}AM$ Geometrically}
\label{sec:derive-similarity}

Now we arrive at the key question: given a linear transformation $T$
represented by matrix $A$ in basis $v$, what is its matrix $B$ in a new
basis $v'$?

Here's the geometric derivation:

\begin{enumerate}
    \item \textbf{Start with a vector in new basis:}
          $\mathbf{x}_{v'}$ is the input (coordinates in $v'$).

    \item \textbf{Convert to old basis (where we know $A$):}
          $\mathbf{x}_v = M_{v' \to v}\,\mathbf{x}_{v'}$

    \item \textbf{Apply the transformation (in old basis):}
          $\mathbf{y}_v = A\,\mathbf{x}_v = A M_{v' \to v}\,\mathbf{x}_{v'}$

    \item \textbf{Convert result back to new basis:}
          $\mathbf{y}_{v'} = M_{v' \to v}^{-1}\,\mathbf{y}_v
                           = M_{v' \to v}^{-1} A M_{v' \to v}\,\mathbf{x}_{v'}$
\end{enumerate}

But by definition, $\mathbf{y}_{v'} = B\,\mathbf{x}_{v'}$ (this is what
we're looking for). Therefore:

\[
    \boxed{\,B = M^{-1} A M\,}
\]

This is the \textbf{similarity transformation}. The same $M$ appears on
both sides because we are transforming within a single vector space —
the domain and codomain are the same space, so they use the same basis
change.

% ---- B.5 ----
\subsection{Same Space vs. Different Spaces: The Key Distinction}
\label{sec:same-vs-different}

The derivation above assumed \textbf{domain = codomain}. This is why the
same matrix $M$ appears twice:
\[
    \underbrace{\mathbf{y}}_{\text{same space}} = A\underbrace{\mathbf{x}}_{\text{same space}}
\]

\bigskip
\noindent\textbf{What happens when domain $\neq$ codomain?}

Consider the impedance matrix $Z$: it maps \textbf{currents} (one space)
to \textbf{voltages} (a different space). These are fundamentally
different physical quantities, governed by different laws:
\begin{itemize}
    \item Currents obey KCL (sum at a node)
    \item Voltages obey KVL (sum around a loop)
\end{itemize}

When we change basis, we must transform each space independently:
\begin{itemize}
    \item $T^i$: transforms current coordinates (from per-line to DM--CM)
    \item $T^v$: transforms voltage coordinates (from per-line to DM--CM)
\end{itemize}

The transformation becomes:
\[
    Z_{\text{new}} = (T^v)^{-1} \, Z \, (T^i)^{-1}
\]

\bigskip
\noindent\textbf{Why this is not similarity transformation:}

\begin{itemize}
    \item Similarity: $A_{\text{new}} = P^{-1}AP$ — same $P$ on both sides
    \item Circuit:   $Z_{\text{new}} = T^v Z(T^i)^{-1}$ — different matrices!
\end{itemize}

The circuit uses \textbf{two different} transformations because the
domain (currents) and codomain (voltages) are different vector spaces.

% ---- B.6 ----
\subsection{Numerical Example: Similarity Transformation ($A: V \to V$)}
\label{sec:numerical-similarity}

For a linear map within a single space, use the same $P$ on both sides.

\bigskip
\noindent\textbf{Setup:} $V = \mathbb{R}^2$, standard basis.
\[
    A = \begin{bmatrix} 3 & 1 \\ 0 & 2 \end{bmatrix}, \qquad
    P = \begin{bmatrix} 1 & 1 \\ 0 & 1 \end{bmatrix}, \qquad
    P^{-1} = \begin{bmatrix} 1 & -1 \\ 0 & 1 \end{bmatrix}.
\]

\noindent\textbf{Similarity transform (same $P$ on both sides):}
\[
    A_{\text{new}} = P^{-1} A P
    = \begin{bmatrix} 1 & -1 \\ 0 & 1 \end{bmatrix}
      \begin{bmatrix} 3 & 1 \\ 0 & 2 \end{bmatrix}
      \begin{bmatrix} 1 & 1 \\ 0 & 1 \end{bmatrix}
    = \begin{bmatrix} 3 & 2 \\ 0 & 2 \end{bmatrix}.
\]

\noindent\textbf{Key point:} Same $P$ on both sides because domain and
codomain are the same space $\mathbb{R}^2$.

% ---- B.7 ----
\subsection{Numerical Example: Circuit Transformation ($Z: V_i \to V_v$)}
\label{sec:numerical-circuit}

For impedance, domain $\neq$ codomain, so we need two different matrices.
We use $N = 2$ lines with the paper's actual definitions.

\bigskip
\noindent\textbf{Setup:}
\begin{itemize}
    \item Impedance matrix (line 1: $2\,\Omega$, line 2: $6\,\Omega$, asymmetric):
        \[
            Z = \begin{bmatrix} 2 & 0 \\ 0 & 6 \end{bmatrix}\ \Omega
        \]
    \item Current transformation $T^i$ (KCL: half-differences and sum):
        \[
            T^i = \begin{bmatrix} \tfrac{1}{2} & -\tfrac{1}{2} \\[4pt] 1 & 1 \end{bmatrix},
            \qquad
            (T^i)^{-1} = \begin{bmatrix} 1 & \tfrac{1}{2} \\[4pt] -1 & \tfrac{1}{2} \end{bmatrix}
        \]
    \item Voltage transformation $T^v$ (KVL: full difference and average):
        \[
            T^v = \begin{bmatrix} 1 & -1 \\[4pt] \tfrac{1}{2} & \tfrac{1}{2} \end{bmatrix}
        \]
\end{itemize}

\bigskip
\noindent\textbf{Step 1 — compute $Z\,(T^i)^{-1}$:}
\[
    Z\,(T^i)^{-1}
    = \begin{bmatrix} 2 & 0 \\ 0 & 6 \end{bmatrix}
      \begin{bmatrix} 1 & \tfrac{1}{2} \\ -1 & \tfrac{1}{2} \end{bmatrix}
    = \begin{bmatrix} 2 & 1 \\ -6 & 3 \end{bmatrix}
\]

\bigskip
\noindent\textbf{Step 2 — apply $T^v$ on the left:}
\[
    Z_{\text{DCM}} = T^v \, Z \, (T^i)^{-1}
    = \begin{bmatrix} 1 & -1 \\ \tfrac{1}{2} & \tfrac{1}{2} \end{bmatrix}
      \begin{bmatrix} 2 & 1 \\ -6 & 3 \end{bmatrix}
    = \begin{bmatrix} 8 & -2 \\ -2 & 2 \end{bmatrix}\ \Omega
\]

\noindent\textbf{Reading the result:}
\begin{itemize}
    \item Diagonal entries: DM impedance $= 8\,\Omega$, CM impedance $= 2\,\Omega$.
    \item Off-diagonal entry $-2\,\Omega \neq 0$ — DM and CM are coupled because
          $Z_1 \neq Z_2$ (asymmetric circuit).
\end{itemize}

\bigskip
\noindent\textbf{Sanity check — symmetric case ($Z_1 = Z_2 = 2\,\Omega$):}
\[
    Z_{\text{DCM}}^{\text{sym}}
    = \begin{bmatrix} 1 & -1 \\ \tfrac{1}{2} & \tfrac{1}{2} \end{bmatrix}
      \begin{bmatrix} 2 & 0 \\ 0 & 2 \end{bmatrix}
      \begin{bmatrix} 1 & \tfrac{1}{2} \\ -1 & \tfrac{1}{2} \end{bmatrix}
    = \begin{bmatrix} 4 & 0 \\ 0 & 1 \end{bmatrix}\ \Omega \quad \checkmark
\]
Diagonal — no coupling when the circuit is symmetric.

\noindent\textbf{Key point:} Different matrices $T^i$ and $T^v$ because
currents and voltages live in different physical spaces (KCL vs.\ KVL).

% ---- B.8 ----
\subsection{Summary}
\label{sec:summary}

\begin{table}[h!]
    \centering
    \begin{tabular}{|l|c|c|}
        \hline
        & Similarity Transform & Circuit Transform \\
        \hline
        \textbf{Map type} & $A: V \to V$ & $Z: V_i \to V_v$ \\
        \hline
        \textbf{Spaces} & Domain = Codomain & Domain $\neq$ Codomain \\
        \hline
        \textbf{Matrices} & Same $P$ on both sides & $T^i$ for currents, $T^v$ for voltages \\
        \hline
        \textbf{Formula} & $A_{\text{new}} = P^{-1}AP$ & $Z_{\text{new}} = T^v Z(T^i)^{-1}$ \\
        \hline
        \textbf{Physical meaning} & Same transformation, new basis & Different physical quantities \\
        \hline
    \end{tabular}
\end{table}

The geometric interpretation clarifies why: similarity transformation
keeps the same vector space (just re-coordinates it), while the circuit
transformation bridges two different physical spaces (currents to
voltages), requiring two independent basis changes.
